\documentclass[11pt,a4paper]{article}
\usepackage[a4paper,margin=2.05cm]{geometry}
\usepackage[spanish,es-noquoting]{babel}
\usepackage[T1]{fontenc}
\usepackage[utf8]{inputenc}
\usepackage{lmodern,microtype}
\usepackage{amsmath,amssymb,amsthm,mathtools}
\usepackage{booktabs,array,longtable}
\usepackage{xcolor}
\usepackage{hyperref}
\hypersetup{colorlinks=true,linkcolor=blue!50!black,urlcolor=blue!50!black}
\setlength{\parindent}{0pt}
\setlength{\parskip}{0.62em}
\newtheorem{teorema}{Teorema}
\newtheorem{proposicion}[teorema]{Proposición}
\newtheorem{corolario}[teorema]{Corolario}
\newtheorem{definicion}[teorema]{Definición}
\newtheorem{observacion}[teorema]{Observación}
\newcommand{\dr}{\operatorname{dr}_9}
\newcommand{\APP}{\mathrm{APP}}
\newcommand{\F}{\mathbb F}
\newcommand{\Z}{\mathbb Z}
\newcommand{\Tr}{\operatorname{tr}}
\newcommand{\Span}{\operatorname{span}}
\title{\bfseries El núcleo aritmético de APP\\
\large Radical $3$--$6$--$9$, compresión holográfica, bloque central y elevación ángulo--contraángulo}
\author{Ventana técnica HMT}
\date{}
\begin{document}
\maketitle

\begin{abstract}
Se estudia la hoja multiplicativa de la célula APP de orden nueve como tabla regular del anillo local $\Z/9\Z$. La red formada por las filas y columnas $3,6,9$ se identifica con el ideal maximal nilpotente $(3)$, cuyo cuadrado es cero. Al reagrupar la tabla por clases módulo tres, la matriz $9\times9$ se comprime exactamente en una matriz $3\times3$ de medias enteras y firma lorentziana. El bloque central $\bigl(\begin{smallmatrix}7&2\\2&7\end{smallmatrix}\bigr)$ se descompone en canales two-way de autovalores $9$ y $5$, y sus lecturas decimales $72,27,77,22$ codifican invariantes globales de la red radical: $45$, $54$, $55$ y $99$. Se formula finalmente la elevación física canónica $7I+2R\mapsto AI+C^\ast R$, que produce los canales $A\pm C^\ast$ y los exponentes $q_\pm$ de Mecánica Dimensional.
\end{abstract}

\section{La hoja multiplicativa como objeto algebraico}

Sea
\[
R=\Z/9\Z,
\]
y representemos el residuo cero por el dígito $9$. Definimos
\[
\dr(n)=1+((n-1)\bmod 9)\in\{1,\dots,9\}.
\]
La hoja multiplicativa de APP es la matriz
\[
M=(M_{ij})_{1\le i,j\le9},\qquad M_{ij}=\dr(ij).
\]
Explícitamente,
\[
M=
\begin{pmatrix}
1&2&3&4&5&6&7&8&9\\
2&4&6&8&1&3&5&7&9\\
3&6&9&3&6&9&3&6&9\\
4&8&3&7&2&6&1&5&9\\
5&1&6&2&7&3&8&4&9\\
6&3&9&6&3&9&6&3&9\\
7&5&3&1&8&6&4&2&9\\
8&7&6&5&4&3&2&1&9\\
9&9&9&9&9&9&9&9&9
\end{pmatrix}.
\]

El anillo $R$ posee grupo de unidades
\[
U=R^\times=\{1,2,4,5,7,8\}
\]
y único ideal maximal
\[
\mathfrak m=(3)=\{0,3,6\},
\]
que en la notación digital es $\{9,3,6\}$. Además,
\[
\mathfrak m^2=0.
\]
Esta igualdad expresa que todo producto de dos elementos de la red $3$--$6$--$9$ colapsa en el cero del anillo, representado por $9$.

\begin{teorema}[Esqueleto radical 3--6--9]
La red formada por las filas y columnas $3,6,9$ es un subobjeto canónico e invariante de APP. Satisface
\[
R\mathfrak m\subseteq\mathfrak m,
\qquad
\mathfrak m R\subseteq\mathfrak m,
\qquad
\mathfrak m^2=0.
\]
Por tanto, separa la dinámica reversible de las unidades de la dinámica nilpotente de cierre.
\end{teorema}

\section{Tres atlas canónicos de la misma matriz}

La matriz $M$ admite tres ordenaciones especialmente informativas.

\subsection*{Atlas decimal}
El orden $1,2,\dots,9$ conserva las semillas visibles, los pares centrales $4,5$ y las lecturas decimales $72,27,77,22$.

\subsection*{Atlas cíclico de unidades}
Ordenando las unidades como potencias del generador $2$,
\[
1,2,4,8,7,5,
\]
el bloque $U\times U$ se convierte en la tabla circulante del grupo cíclico $C_6$:
\[
\begin{pmatrix}
1&2&4&8&7&5\\
2&4&8&7&5&1\\
4&8&7&5&1&2\\
8&7&5&1&2&4\\
7&5&1&2&4&8\\
5&1&2&4&8&7
\end{pmatrix}.
\]
Así, las seis filas reversibles son traslaciones cíclicas de una misma órbita.

\subsection*{Atlas Hensel por clases módulo tres}
Definimos
\[
C_1=\{1,4,7\},\qquad
C_2=\{2,5,8\},\qquad
C_0=\{3,6,9\}.
\]
El conjunto $C_0$ es el ideal nilpotente, mientras $C_1$ y $C_2$ son las dos clases de unidades del cociente $R/\mathfrak m\cong\F_3$.

\section{Compresión holográfica $9\times9\to3\times3$}

Particionamos $M$ en nueve bloques $3\times3$ según $C_1,C_2,C_0$. Sea $S_\times$ la matriz de sumas de bloque. Un cálculo exacto da
\[
\boxed{
S_\times=
\begin{pmatrix}
36&45&54\\
45&36&54\\
54&54&81
\end{pmatrix}.}
\]
Cada bloque contiene nueve entradas, por lo que su matriz de medias es
\[
\boxed{
Q_\times=\frac1{9}S_\times=
\begin{pmatrix}
4&5&6\\
5&4&6\\
6&6&9
\end{pmatrix}.}
\]

\begin{teorema}[Compresión exacta de APP]
La matriz $Q_\times$ conserva la estructura de clases módulo tres, el total de la hoja y la separación entre canales unitarios y nilpotentes. En particular,
\[
\sum_{a,b}(Q_\times)_{ab}=51,
\qquad
\sum_{i,j}M_{ij}=9\cdot51=459.
\]
Además,
\[
\Tr(Q_\times)=17,
\qquad
459=27\cdot17.
\]
\end{teorema}

La hoja aditiva, tomando $A_{ij}=\dr(i+j-1)$, produce de modo análogo
\[
Q_+=
\begin{pmatrix}
4&5&6\\
5&6&4\\
6&4&5
\end{pmatrix},
\]
con
\[
\sum Q_+=45,
\qquad
\Tr(Q_+)=15,
\qquad
\sum_{i,j}A_{ij}=27\cdot15=405.
\]
Por tanto,
\[
Q_\times-Q_+=
\begin{pmatrix}
0&0&0\\
0&-2&2\\
0&2&4
\end{pmatrix},
\]
y se obtiene la jerarquía exacta
\[
\boxed{
\Tr(Q_\times-Q_+)=2,
\quad
\sum(Q_\times-Q_+)=6,
\quad
459-405=27\cdot2=54.}
\]
El defecto global $54$ es una amplificación triádica del defecto comprimido $2$.

\begin{teorema}[Localización del defecto APP]
Las seis filas unitarias de $M$ son permutaciones de $1,\dots,9$ y suman $45$. Las filas $3$ y $6$ suman $54$, mientras la fila $9$ suma $81$. Por tanto,
\[
\Delta_{\APP}
=2(54-45)+(81-45)=54.
\]
Todo el defecto aditivo--multiplicativo se concentra en el sector no unitario $\mathfrak m$.
\end{teorema}

\section{Firma lorentziana y operador de Pell}

La forma $Q_\times$ posee determinante y espectro
\[
\det Q_\times=-9,
\]
\[
\operatorname{Spec}(Q_\times)
=
\{-1,\,9-6\sqrt2,\,9+6\sqrt2\}.
\]
Por tanto tiene firma $(2,1)$.

Sea
\[
v_-=(1,-1,0),\qquad
v_+=(1,1,0),\qquad
v_0=(0,0,1).
\]
Entonces
\[
Q_\times v_-=-v_-.
\]
En el subespacio $\Span\{v_+,v_0\}$, la matriz de $Q_\times$ es
\[
\begin{pmatrix}
9&6\\
12&9
\end{pmatrix}
=3
\begin{pmatrix}
3&2\\
4&3
\end{pmatrix}.
\]
La matriz
\[
H_2=
\begin{pmatrix}
3&2\\
4&3
\end{pmatrix}
\]
pertenece a $SL(2,\Z)$ y tiene autovalores
\[
3\pm2\sqrt2=(1\pm\sqrt2)^2.
\]

\begin{corolario}[Geometría interna del cociente trítico]
La compresión de APP contiene un canal contraorientado de autovalor $-1$ y un canal hiperbólico de Pell gobernado por la unidad $3+2\sqrt2$. La firma lorentziana surge de la propia aritmética módulo nueve.
\end{corolario}

\section{El bloque central $7,2;2,7$}

Las posiciones centrales $4,5$ son opuestas módulo nueve. Su bloque multiplicativo es
\[
\boxed{
B=
\begin{pmatrix}
7&2\\
2&7
\end{pmatrix}.}
\]
Sea
\[
R=
\begin{pmatrix}
0&1\\
1&0
\end{pmatrix},
\qquad R^2=I.
\]
Entonces
\[
B=7I+2R.
\]
Con los proyectores
\[
P_+=\frac{I+R}{2},
\qquad
P_-=\frac{I-R}{2},
\]
se obtiene
\[
\boxed{B=9P_+ +5P_-.}
\]
Los dos canales propios son, por tanto, el cierre nonádico $9$ y el canal pentádico $5$.

Además,
\[
\det B=45,
\qquad
\Tr B=14,
\qquad
\sum B_{ij}=18.
\]
La normalización por $\sqrt{45}$ produce
\[
\frac{B}{\sqrt{45}}=
\exp(\eta_0R),
\qquad
\eta_0=\operatorname{artanh}\frac27
=\frac12\log\frac95.
\]

\section{Lecturas decimales y doble partición de $99$}

La lectura por filas produce
\[
72,\qquad27.
\]
La lectura por diagonales produce
\[
77,\qquad22.
\]
Estas cuatro cantidades satisfacen
\[
\boxed{72+27=77+22=99,}
\]
\[
\boxed{72-27=45=\det B,}
\]
\[
\boxed{77-22=55=\Delta_{\APP}+1,}
\]
\[
\boxed{99=45+54.}
\]
Además,
\[
72=8\cdot9,
\qquad
27=3\cdot9,
\]
por lo que
\[
99=(8+3)\cdot9=11\cdot9.
\]
Y
\[
77=7\cdot11,
\qquad
22=2\cdot11,
\]
por lo que
\[
99=(7+2)\cdot11=9\cdot11.
\]

\begin{teorema}[Doble partición del núcleo]
El bloque central codifica dos particiones duales del mismo cierre $99$:
\[
11=8+3
\quad\text{en la lectura por filas},
\]
\[
9=7+2
\quad\text{en la lectura diagonal}.
\]
La primera separa sector interno octádico y sector triádico; la segunda separa los dos coeficientes del operador two-way.
\end{teorema}

La proyección
\[
\frac{22}{7}
\]
es la lectura racional clásica de $\pi$ asociada al mismo bloque. Debe distinguirse de los invariantes matriciales: es una proyección decimal canónica del atlas natural.

\section{Teorema local--global de la red radical}

La unión de las filas y columnas $3,6,9$ tiene suma total
\[
297.
\]
Su núcleo $C_0\times C_0$ contiene nueve entradas iguales a $9$ y suma
\[
81=3\cdot27.
\]
Los cuatro brazos restantes suman
\[
216=8\cdot27.
\]
Por tanto,
\[
\boxed{297=216+81=(8+3)\cdot27=11\cdot27.}
\]
Pero las lecturas centrales satisfacen
\[
72=8\cdot9,
\qquad
27=3\cdot9.
\]
Luego
\[
\boxed{3\cdot72=216,\qquad 3\cdot27=81,\qquad3\cdot99=297.}
\]

\begin{teorema}[Holograma central de la red $3$--$6$--$9$]
La lectura por filas del bloque central es una compresión exacta, a escala tres, de la descomposición global de la red radical:
\[
72\longleftrightarrow216\quad\text{(brazos)},
\]
\[
27\longleftrightarrow81\quad\text{(núcleo)},
\]
\[
99\longleftrightarrow297\quad\text{(red completa)}.
\]
El núcleo local $7,2;2,7$ contiene así la contabilidad global de la red $3,6,9$.
\end{teorema}

La parte complementaria $U\times U$ suma
\[
162=6\cdot27,
\]
y por tanto
\[
459=(11+6)\cdot27=17\cdot27.
\]

\section{Del bloque central al ángulo y contraángulo}

La interfaz angular de Mecánica Dimensional conserva el álgebra two-way $\Span\{I,R\}$. El levantamiento físico es
\[
\boxed{
B=7I+2R
\quad\longmapsto\quad
B_{A,C^\ast}=AI+C^\ast R.}
\]
Aquí
\[
A=1000\alpha,
\qquad
C^\ast=2(\bar\hbar+\alpha).
\]
La descomposición espectral es
\[
\boxed{
B_{A,C^\ast}
=(A+C^\ast)P_+
+(A-C^\ast)P_- .}
\]
Por tanto, los canales físicos son
\[
\lambda_+=A+C^\ast,
\qquad
\lambda_-=A-C^\ast.
\]
Los canales analíticos del corpus son
\[
\boxed{
q_\pm=
\exp\left[-\frac{\pi}{180}(A\pm C^\ast)\right].}
\]
El producto y el cociente reconstruyen exactamente las coordenadas comunes y diferenciales:
\[
\boxed{
A=-\frac{90}{\pi}\log(q_+q_-),}
\]
\[
\boxed{
C^\ast=\frac{90}{\pi}\log\frac{q_-}{q_+}.}
\]
La conjugación de hoja actúa como
\[
C^\ast\mapsto-C^\ast,
\qquad
q_+\leftrightarrow q_-.
\]

El operador físico normalizado es
\[
\frac{B_{A,C^\ast}}{\sqrt{A^2-(C^\ast)^2}}
=
\exp\bigl(\eta_{A,C^\ast}R\bigr),
\]
con
\[
\eta_{A,C^\ast}
=
\operatorname{artanh}\frac{C^\ast}{A}
=
\frac12\log\frac{A+C^\ast}{A-C^\ast}.
\]
Así, la estructura hiperbólica del bloque entero se conserva en su elevación física.

\section{Del canal local al residuo dodecafásico}

El bloque central local y el bloque unitario comprimido son
\[
B=\begin{pmatrix}7&2\\2&7\end{pmatrix}
=9P_+ +5P_-,
\]
\[
C=\begin{pmatrix}4&5\\5&4\end{pmatrix}
=9P_+ -P_-.
\]
Por tanto,
\[
\boxed{B-C=6P_-.}
\]
La compresión conserva el canal común $9$ y desplaza exclusivamente el canal contraorientado en seis unidades:
\[
5\longmapsto-1.
\]
Si la rueda dodecafásica normaliza el canal contraorientado por doce fases, aparece el residuo
\[
\boxed{-\frac1{12}.}
\]
Esta igualdad identifica el menos un doceavo como sombra dodecafásica del autovalor contraorientado $-1$ del núcleo comprimido.

\section{Conclusión}

La hoja multiplicativa APP es simultáneamente:
\begin{itemize}
\item la tabla regular del anillo local $\Z/9\Z$;
\item una dinámica reversible sobre el grupo cíclico de seis unidades;
\item una extensión nilpotente gobernada por la red $3,6,9$;
\item una forma comprimida de firma lorentziana;
\item un operador central two-way con canales $9$ y $5$;
\item un atlas decimal que produce $72,27,77,22$ y $22/7$;
\item el esqueleto algebraico de la interfaz física $A,C^\ast$.
\end{itemize}

La red radical y el bloque central están unidos por identidades exactas de escala. La matriz $9\times9$ no contiene una colección desordenada de productos: contiene tres atlas de una misma aritmética compactada. El atlas decimal revela semillas; el atlas cíclico revela órbitas; el atlas Hensel revela radical, firma y memoria. Esta pluralidad controlada de lecturas constituye una formulación precisa de aritmética holográfica.

\end{document}
